% !TEX root = main.tex
\clearpage
\section{Code guide: data pipeline v0}

Every pipeline stage reads one JSON object per line from standard input and
writes surviving documents to standard output. This Unix-filter design makes
the stages composable and allows rejected or intermediate rows to be inspected.

\subsection{Canonical document flow}

\small\ttfamily
\begin{longtable}{@{}p{0.22\textwidth}p{0.64\textwidth}@{}}
  ingest & source-specific API or archive $\rightarrow$ id, source, url, text \\
  langid.py & add lang, script, code\_mix \\
  filters.py & remove low-quality prose; optionally record reject reasons \\
  dedup\_minhash.py & remove near-duplicate documents \\
  decontaminate.py & remove documents overlapping evaluation text \\
  mix.py & calculate domain sampling rates and write a snapshot manifest \\
  prepare\_data.py & tokenize text and write train.bin, val.bin, meta.json
\end{longtable}
\normalsize\rmfamily

Code bypasses prose quality filters because source code legitimately violates
their assumptions about word length, symbols, bullets, and stopwords.

\subsection{Legal register}

\texttt{legal\_register.md} is an append-only compliance ledger for
training-data sources. Each row records:
\begin{itemize}
    \setlength{\itemsep}{0pt}
  \item source and content type;
  \item legal basis and license/terms;
  \item reviewer and review date;
  \item restrictions or outstanding compliance work; and
  \item a stable source ID such as \texttt{S010}.
\end{itemize}
Ingestion scripts attach that ID to every generated document, making data
provenance traceable. The admin app can add review-pending entries, and the
register is copied to object storage alongside corpus snapshots.

The ingestion scripts require a \texttt{--source-id}, but do not verify that it
exists in the register or has compliance approval.

\subsection{Ingestion adapters}

\texttt{ingest\_hf\_rows.py} paginates the Hugging Face datasets-server API in
100-row requests. It records a stable row-based ID and the legal-register
source ID, backs off on HTTP 429 responses, and skips empty text.

\texttt{ingest\_wikipedia.py} supports curated title lists and random pages. It
fetches one complete extract per request, follows redirects, removes stubs
under 300 characters, and deduplicates titles within a run.

\texttt{ingest\_github\_code.py} downloads tarballs only for four explicitly
allowlisted repositories. It emits one document per Python file, records the
repository license, and skips files larger than 200KB.

\texttt{run\_ingest.sh} orchestrates the six domain pools and then applies the
cleaning chain. Its current sources are useful for mechanics but are not a
production corpus: Wikipedia topic lists are narrow, the GitHub allowlist is
tiny, and the Hugging Face rows API is a laptop-scale sampler.

\subsection{Language and code-mix identification}

\texttt{script\_profile()} counts alphabetic characters within explicit Unicode
ranges. Native-script text is mapped coarsely to a language; more than 15\%
Latin letters marks it as code-mixed. Latin-dominant text becomes Hinglish when
more than 5\% of whitespace tokens match a small marker lexicon.

This approach is transparent, but it cannot distinguish Hindi from Marathi in
Devanagari, Bengali from Assamese in Bengali script, or English containing a
few Hindi discourse markers from genuine Hinglish. It is a routing heuristic,
not a production language identifier.

\subsubsection{Mixing and snapshot identity}

\texttt{mix.py} does not actually concatenate or tokenize the data yet. Its
job is to create a mixing plan / manifest that says: for the final training
corpus, how much should come from each domain, and how many times do I need
to sample each pool to hit that target?

Suppose you want a 100M-token training corpus. Then \texttt{mix.py}
calculates each domain's target as $\text{weight}\times 100\text{M}$:
\begin{center}
  \renewcommand{\arraystretch}{1.15}
  \small
  \begin{tabular}{@{}l r l@{}}
    \textbf{Domain} & \textbf{Weight} & \textbf{Target tokens} \\
    \hline
    \texttt{global\_english} & 50\% &
    $0.50\times 100\text{M}=50\text{M}$ \\
    \texttt{indian\_english} & 17.5\% &
    $0.175\times 100\text{M}=17.5\text{M}$ \\
    \texttt{indian\_languages} & 12.5\% &
    $0.125\times 100\text{M}=12.5\text{M}$ \\
    \texttt{code} & 10\% &
    $0.10\times 100\text{M}=10\text{M}$ \\
    \texttt{finance} & 7.5\% &
    $0.075\times 100\text{M}=7.5\text{M}$ \\
    \texttt{math} & 2.5\% &
    $0.025\times 100\text{M}=2.5\text{M}$ \\
  \end{tabular}
\end{center}
Then it asks: how much clean data do I actually have in each pool?

\begin{paragraphblock}{The default weights}
\begin{lstlisting}[style=shadedpython,breaklines=true]
DEFAULT_WEIGHTS = {
    "global_english": 0.50,
    "indian_english": 0.175,
    "indian_languages": 0.125,
    "code": 0.10,
    "finance_econ_law": 0.075,
    "math_reasoning": 0.025,
}
\end{lstlisting}
  These sum to 100\%.
\end{paragraphblock}

\begin{paragraphblock}{Input config}
\begin{lstlisting}[style=shadedpython,breaklines=true]
python mix.py \
    --config mix_v1.json \
    --total-tokens 100e9 \
    --out snapshots/
\end{lstlisting}
  The config is expected to look conceptually like:
\begin{lstlisting}[style=shadedpython,breaklines=true]
{
  "global_english": {
    "paths": ["pools/global_english/clean.jsonl"],
    "tokens": 40000000000,
    "weight": 0.50
  },
  "indian_languages": {
    "paths": ["pools/indian_languages/clean.jsonl"],
    "tokens": 5000000000,
    "weight": 0.125
  }
}
\end{lstlisting}
  For each domain it needs \texttt{paths} (where cleaned data lives),
  \texttt{tokens} (how many clean tokens exist in that pool), and
  \texttt{weight} (desired fraction of the final training corpus). If
  \texttt{weight} is absent, it uses \texttt{DEFAULT\_WEIGHTS}.
\end{paragraphblock}

\begin{paragraphblock}{Validate weights}
\begin{lstlisting}[style=shadedpython,breaklines=true]
wsum = sum(
    d.get("weight", DEFAULT_WEIGHTS.get(name, 0))
    for name, d in cfg.items()
)
if abs(wsum - 1.0) > 1e-6:
    sys.exit(...)
\end{lstlisting}
  makes sure they equal 1 (100\%).
\end{paragraphblock}

\begin{paragraphblock}{Calculate each domain's target}
  For each domain, \texttt{target = w * total} means target tokens $=$ domain
  weight $\times$ total desired tokens. Example: total $=100\text{M}$, global
  English weight $=0.50$ $\Rightarrow$ target $=50\text{M}$ tokens.
\end{paragraphblock}

\begin{paragraphblock}{Look at available pool size}
  \texttt{pool = float(d["tokens"])}. Suppose the cleaned English pool
  contains 40M tokens but the target is 50M --- you don't have enough unique
  English data to provide 50M tokens once. So the script calculates how many
  epochs over that pool are needed.
\end{paragraphblock}

\begin{paragraphblock}{What does epochs mean here?}
  \texttt{epochs = target / pool}. If target $=50\text{M}$ and pool
  $=40\text{M}$, then epochs $=1.25$: read the full pool once, plus another
  25\% of it.
  \begin{itemize}
      \setlength{\itemsep}{0pt}
    \item \texttt{epochs = 1.0} --- use the pool once
    \item \texttt{epochs = 2.0} --- repeat the whole pool twice
    \item \texttt{epochs = 0.5} --- sample about half the pool
    \item \texttt{epochs = 3.5} --- use the pool about $3\frac{1}{2}$ times
  \end{itemize}
  This is not model-training epochs yet. It means how much of that data pool
  needs to be sampled into the intended corpus.
\end{paragraphblock}

\begin{paragraphblock}{Why cap at 4 epochs?}
\begin{lstlisting}[style=shadedpython]
MAX_EPOCHS = 4.0
capped = min(epochs, MAX_EPOCHS)
\end{lstlisting}
  Suppose Indian language pool $=2\text{M}$ tokens and target $=12.5\text{M}$:
  required epochs $=6.25$, but the script refuses to go beyond $4\times$
  because the project has chosen not to over-repeat small datasets
  excessively.
\end{paragraphblock}

\begin{paragraphblock}{Calculate what you can actually get}
  \texttt{got = capped * pool}. Using the same example: pool $=2\text{M}$,
  capped epochs $=4$ $\Rightarrow$ \texttt{got}$=8\text{M}$ tokens, but target
  was $12.5\text{M}$ --- so there's a shortage.
\end{paragraphblock}

\begin{paragraphblock}{Shortfall}
  \texttt{shortfall\_tokens = max(0.0, target - got)}. Example: target
  $=12.5\text{M}$, got $=8\text{M}$ $\Rightarrow$ shortfall $=4.5\text{M}$
  tokens. The manifest records that explicitly.
\end{paragraphblock}

\begin{paragraphblock}{What gets stored per domain?}
\begin{lstlisting}[style=shadedpython,breaklines=true]
domains[name] = {
    "paths": d["paths"],
    "pool_tokens": pool,
    "weight": w,
    "target_tokens": target,
    "epochs": round(capped, 4),
    "sampled_tokens": round(got),
    "shortfall_tokens": round(max(0.0, target - got)),
}
\end{lstlisting}
  creates something conceptually like:
\begin{lstlisting}[style=shadedpython,breaklines=true]
{
  "global_english": {
    "paths": ["pools/global_english/clean.jsonl"],
    "pool_tokens": 40000000,
    "weight": 0.50,
    "target_tokens": 50000000,
    "epochs": 1.25,
    "sampled_tokens": 50000000,
    "shortfall_tokens": 0
  }
}
\end{lstlisting}
\end{paragraphblock}

\begin{paragraphblock}{Build the manifest}
\begin{lstlisting}[style=shadedpython,breaklines=true]
manifest = {
    "total_tokens_requested": total,
    "max_epochs": MAX_EPOCHS,
    "domains": domains,
}
\end{lstlisting}
  This is the recipe for the training corpus.
\end{paragraphblock}

\begin{paragraphblock}{Generate a snapshot ID}
\begin{lstlisting}[style=shadedpython,breaklines=true]
blob = json.dumps(
    manifest, sort_keys=True, ensure_ascii=False
).encode()
hashlib.sha256(blob)
\end{lstlisting}
  computes a SHA-256 fingerprint. Only the first 16 hex characters are used:
\begin{lstlisting}[style=shadedpython,language={},basicstyle=\scriptsize\ttfamily]
snap_id = "snap-" + ...[:16]
# e.g. snap-6a97923395c5573e
\end{lstlisting}
  Why hash the manifest? The snapshot ID identifies the exact corpus recipe.
  Change global English from 50\% to 45\% and the SHA-256 changes --- so
  experiments stay traceable: ``Model run \#37 used
  \texttt{snap-6a97923395c5573e}''.
\end{paragraphblock}

\begin{paragraphblock}{Snapshots are immutable}
\begin{lstlisting}[style=shadedpython,breaklines=true]
if os.path.exists(path):
    sys.exit(
        f"{path} exists — snapshots are immutable, never overwrite"
    )
\end{lstlisting}
  Once a snapshot exists, don't modify it. If you change the data mix, create
  a \emph{new} snapshot instead of rewriting the old one.
\end{paragraphblock}

\begin{paragraphblock}{Save snapshot}
  Creates \texttt{snapshots/snap-6a97923395c5573e.json}. Important: this file
  is a \emph{manifest}, not the actual 100B tokens of training text. It says
  how to construct/sample the corpus.
\end{paragraphblock}

\begin{paragraphblock}{Prints a summary}
  For each domain it prints something like:
\begin{lstlisting}[style=shadedpython,language={},basicstyle=\scriptsize\ttfamily]
global_english     w=0.500 epochs=0.72 sampled=5e+10
indian_english     w=0.175 epochs=1.35 sampled=1.75e+10
indian_languages   w=0.125 epochs=4.00 sampled=8e+09 SHORTFALL 4.5e+09
\end{lstlisting}
  This lets you immediately see target weight, required sampling/repetition,
  actual sampled tokens, and shortfalls.
\end{paragraphblock}

\begin{paragraphblock}{Warning when data is insufficient}
  If any domain needs more than 4 epochs, \texttt{shortfall = True}. At the
  end:
\begin{lstlisting}[style=shadedpython,breaklines=true]
if shortfall:
    print(
        "WARNING: shortfalls present — "
        "get more data or rebalance weights"
    )
\end{lstlisting}
  You then have two choices: collect more data for that domain, or
  change/rebalance the desired weights.
\end{paragraphblock}

\begin{paragraphblock}{Full example}
  Imagine \texttt{TOTAL REQUESTED = 100M} tokens. Then \texttt{mix.py}
  calculates:
  \begin{center}
    \renewcommand{\arraystretch}{1.2}
    \small
    \begin{tabular}{@{}l r r r r r l@{}}
      \textbf{Domain} & \textbf{Weight} & \textbf{Pool} &
      \textbf{Target} & \textbf{Epochs} & \textbf{Sampled} &
      \textbf{Notes} \\
      \hline
      global English & 50\% & 200M & 50M & 0.25 & 50M &
      sample 25\% of pool \\
      Indian English & 17.5\% & 20M & 17.5M & 0.875 & 17.5M & \\
      Indian languages & 12.5\% & 2M & 12.5M & 4.0 & 8M &
      raw $6.25$, capped; shortfall $4.5$M \\
      code & 10\% & 20M & 10M & 0.5 & 10M & \\
      finance & 7.5\% & 10M & 7.5M & 0.75 & 7.5M & \\
      math & 2.5\% & 5M & 2.5M & 0.5 & 2.5M & \\
    \end{tabular}
  \end{center}
  Then it saves \texttt{snap-xxxxxxxxxxxxxxxx.json} containing all those
  decisions.
\end{paragraphblock}

\begin{paragraphblock}{Where \texttt{mix.py} sits in the pipeline}
  End-to-end mental model:
\begin{lstlisting}[style=shadedpython,language={},basicstyle=\scriptsize\ttfamily]
Internet data
    |
    v
ingest
    |
    v
raw.jsonl
    |
    v
clean
    |
    v
clean.jsonl
    |
    +-- tokenizer branch -> tokenizer.json
    |
    +-- model-mix branch -> snapshot.json
                         |
                         v
          tokenizer.json + snapshot.json
                         |
                         v
                  prepare_data.py
                         |
                         v
                    train/val bins
                         |
                         v
                    MODEL TRAINING
\end{lstlisting}

  So \texttt{clean.jsonl} is the common fork point, and
  \texttt{prepare\_data.py} is the join point.

\begin{lstlisting}[style=shadedpython,language={},basicstyle=\scriptsize\ttfamily]
                        DATA SOURCES

         FineWeb / Sangraha / Wikipedia / GitHub
                            |
                            v
                         INGEST
                            |
                            v
                        raw.jsonl
                            |
                            v
                     DATA CLEANING

                         langid.py
              (detect language, script, and code-mixing)
                            |
                            v
                        filters.py
       (remove low-quality / malformed prose using
        length, repetition, symbol, and language checks)
                    (skipped for code)
                            |
                            v
                    dedup_minhash.py
       (remove near-duplicate documents using 5-word
        shingles + MinHash/LSH + Jaccard similarity)
                            |
                            v
                   decontaminate.py
       (remove any training document overlapping
        BachattBench evals via 13-word n-grams)
                            |
                            v

                   CLEANED DATA POOLS
              pools/<domain>/clean.jsonl
                            |
                +-----------+-----------+
                |                       |
                v                       v

        TOKENIZER WORKSTREAM      MODEL DATA WORKSTREAM
        ====================      =====================

          sample_corpus.py                mix.py
                |                           |
        take capped sample          read domain config
        from each domain            containing:
                |                   * clean paths
                v                   * pool token counts
          corpus/*.txt              * target weights
                |                           |
                v                           v
        train_tokenizer.py           calculate target
                |                    tokens per domain
                v                           |
         tokenizer.json                     v
                |                    calculate required
                v                         epochs
        eval_fertility.py                   |
                |                           v
        measure tokens/word          cap repetition at 4x
        across languages                    |
                |                           v
        validate tokenizer           record shortfalls
                |                           |
                |                           v
                |                    snapshot manifest
                |                           |
                |                           v
                |            snap-6a97923395c5573e.json
                |                           |
                +-----------+---------------+
                            |
                            v
                     prepare_data.py
                            |
              +-------------+-------------+
              |                           |
              v                           v

       snapshot manifest             tokenizer.json

       tells prepare_data:           tells prepare_data:
       WHAT data to use              HOW text becomes
       and in what proportions       token IDs

              |                           |
              +-------------+-------------+
                            |
                            v
                  PHYSICALLY SAMPLE DATA
                            |
                            v
                      TOKENIZE TEXT
                            |
                            v
                    TOKEN ID STREAM
                            |
                     train / val split
                            |
                  +---------+---------+
                  v                   v
              train bins           val bins
                  |
                  +---------+---------+
                            v
                      MODEL TRAINING
\end{lstlisting}

  \newpage
  Slightly more detailed version of \texttt{mix.py}:
\begin{lstlisting}[style=shadedpython,language={},basicstyle=\scriptsize\ttfamily]
clean.jsonl pools
       |
       v
     mix.py
       |
       +-- desired domain weights
       |
       |   global English      50%
       |   Indian English      17.5%
       |   Indian languages    12.5%
       |   code                10%
       |   finance/econ/law     7.5%
       |   math/reasoning       2.5%
       |
       +-- available pool sizes
       |
       v
target tokens per domain

target = weight x total_tokens
       |
       v
required epochs

epochs = target_tokens / pool_tokens
       |
       v
cap at MAX_EPOCHS = 4
       |
       +-- if target achievable
       |      -> sampled_tokens = target
       |
       +-- if pool too small
              -> repeat at most 4x
              -> record shortfall
       |
       v
snapshot manifest
       |
       v
snap-xxxxxxxxxxxxxxxx.json
\end{lstlisting}

  The most important distinction: \texttt{mix.py} does not itself build the
  final corpus. It computes and freezes the \emph{recipe} for how the corpus
  should be mixed.

  Compared with \texttt{sample\_corpus.py}:
  \begin{itemize}
      \setlength{\itemsep}{0pt}
    \item \texttt{sample\_corpus.py} builds a balanced \emph{tokenizer}-training
      sample (equal byte cap).
    \item \texttt{mix.py} defines the \emph{model}-training mixture (target
        domain weights, available token counts, repetition limits, versioned
      snapshot).
  \end{itemize}

  Why per-domain rather than per-language at the top level: the mix weights
  in the charter are defined over domains (50\% global English, 12.5\%
  Indian languages, \ldots), so \texttt{mix.py} samples at that granularity
  --- but per-language files are kept inside \texttt{indian\_languages/}
  precisely so tokenizer work can rebalance by language (that's what
  \texttt{sample\_corpus.py} did for the balanced tokenizer corpus).

  \texttt{sample\_corpus.py} walks each domain directory and copies documents
  from \texttt{clean.jsonl} into one text file until a byte cap is reached. The
  cap prevents the largest pool from completely dominating BPE merge counts.

  The implementation is deliberately simple, but two details matter:

  \begin{itemize}
      \setlength{\itemsep}{0pt}
    \item it takes the first documents in file order rather than a
      seeded random sample; and
    \item it applies an equal byte cap, not the target corpus-mix weights.
  \end{itemize}

  Before the tokenizer freeze, sampling should be seeded, content-hashed, and
  stratified by domain, language, source, and document length.
\end{paragraphblock}
